\chapter{Parts i construccions universals}\label{cap:exercicis-parts}

Aquests exercicis proposats corresponen al capítol homònim de la part teòrica (capítol~\ref{cap:parts} de la Part~\ref{part:teoria}) i es resolen amb les tècniques dels exercicis resolts de la part pràctica (capítol~\ref{cap:practica-parts} de la Part~\ref{part:practica}). Cada secció remet al bloc de resolts que en dona el model. Els marcats amb una indicació són els més exigents.

\section{Cancel·lació i composició de mono/epi}

\begin{enumerate}
  \item Demostreu que si \(m:A\rightarrowtail B\) és un monomorfisme i \(u:X\to A\) és qualsevol, aleshores \(m\circ u=m\circ u'\) implica \(u=u'\); enuncieu i demostreu la propietat dual per a un epimorfisme. (Model: secció~\ref{prac:cancel-lacio}.)

  \item Siguin \(m:A\rightarrowtail B\) i \(n:B\rightarrowtail C\) monomorfismes. Demostreu que \(n\circ m\) és un monomorfisme i, dualment, que la composició de dos epimorfismes és un epimorfisme. \begin{indicacio}Descomponeu \((n\circ m)\circ u\) amb l'associativitat i cancel·leu primer \(n\), després \(m\), com a la proposició~\ref{prop:composicio-mono-epi}.\end{indicacio}

  \item Siguin \(f:A\to B\) i \(g:B\to C\). Demostreu que si \(g\circ f\) és un epimorfisme, aleshores \(g\) és un epimorfisme. (Dual de l'exercici resolt de cancel·lació posterior.)

  \item Doneu, en el model clàssic, un exemple de \(f\) i \(g\) amb \(g\circ f\) mono però \(g\) no mono, per veure que la conclusió de l'exercici anterior no es pot enfortir a \enquote{\(g\) és mono}. \begin{indicacio}Preneu \(g\) no injectiva però injectiva sobre la imatge de \(f\).\end{indicacio}
\end{enumerate}

\section{Seccions, retraccions i exhaustivitat generalitzada}

\begin{enumerate}[resume]
  \item Demostreu que tota retracció és un epimorfisme (la meitat de la proposició~\ref{prop:seccio-mono-retraccio-epi} que no s'ha resolt al text). \begin{indicacio}Dualitzeu l'argument de \enquote{tota secció és un monomorfisme}: component per la dreta amb la secció.\end{indicacio}

  \item Sigui \(e:A\to B\) un epimorfisme escindit amb secció \(s\). Demostreu que \(s\) és un monomorfisme i que \(e\circ s=\id_B\) caracteritza \(s\) entre les possibles seccions. 

  \item Comproveu que, en el model clàssic, una aplicació exhaustiva sobre els elements globals no és necessàriament una retracció si no es disposa d'una elecció de preimatges; relacioneu-ho amb l'observació~\ref{obs:exhaustiva-generalitzada-global}. \begin{indicacio}El pas de \enquote{exhaustiva} a \enquote{escindida} requereix triar, per a cada element de \(B\), una preimatge; sense elecció, la secció pot no existir.\end{indicacio}
\end{enumerate}

\section{Parts, inclusió i pertinença}

\begin{enumerate}[resume]
  \item Demostreu que la relació d'inclusió \(\subseteq_A\) és reflexiva i transitiva sobre les parts de \(A\), i que \(m\subseteq_A n\) i \(n\subseteq_A m\) impliquen que \(m\) i \(n\) són parts equivalents. (Model: secció~\ref{prac:parts-inclusio}; per a l'antisimetria vegeu la proposició~\ref{prop:equivalencia-parts}.)

  \item Siguin \(m,n\) parts de \(A\). Demostreu la implicació que falta a la proposició~\ref{prop:inclusio-pertinenca}: si per a tot \(x:X\to A\) es té \(x\in_X m\Rightarrow x\in_X n\), aleshores \(m\subseteq_A n\). \begin{indicacio}Apliqueu la hipòtesi a l'element generalitzat \enquote{genèric} \(x=m:S\to A\), que pertany a \(m\) via \(\id_S\); el testimoni resultant és el \(k\) buscat.\end{indicacio}

  \item Demostreu que si \(x\in_X m\) i \(u:Y\to X\) és qualsevol, aleshores \(x\circ u\in_Y m\). \begin{indicacio}Si \(\bar x\) testimonia \(x\in_X m\), comproveu que \(\bar x\circ u\) testimonia \(x\circ u\in_Y m\).\end{indicacio}
\end{enumerate}

\section{Sumes, característiques i \texorpdfstring{\(\mathbf2\)}{2}}

\begin{enumerate}[resume]
  \item Donades \(f:A\to C\) i \(g:B\to C\), demostreu que l'aplicació definida per casos \([f,g]:A+B\to C\) és única amb les propietats \([f,g]\circ i_A=f\) i \([f,g]\circ i_B=g\). (Model: exercici resolt d'unicitat per casos, secció~\ref{prac:sumes}.)

  \item Demostreu que, si existeixen aplicacions \(f\colon A\to B\) i \(g\colon B\to A\), les dues injeccions \(i_0,i_1\) de la suma \(A+B\) són monomorfismes escindits. \begin{indicacio}Construïu una retracció de cada injecció com al model del cograf (proposició~\ref{prop:cograf-retraccio}); una secció és mono. Que les injeccions siguin sempre monomorfismes, sense aquesta hipòtesi, es veurà al capítol~\ref{cap:valors}.\end{indicacio}

  \item Demostreu que dues parts de \(A\) amb la mateixa aplicació característica \(\chi:A\to\mathbf2\) són equivalents, i recíprocament que parts equivalents tenen la mateixa característica. (Model: secció~\ref{prac:caracteristica}; useu l'axioma~\ref{ax:representacio}.)

  \item Comproveu que la característica de la part total \(\id_A\) és l'aplicació constant \(1\circ{!}_A\), i demostreu que una part \(m\colon S\rightarrowtail A\) amb \(\chi_m=1\circ{!}_A\) és equivalent a la total i, per tant, \(m\) és un isomorfisme. \begin{indicacio}Per a la segona part, \(\id_A\in_A m\) dona una secció de \(m\); combineu-ho amb el fet que \(m\) és mono, com a la proposició~\ref{prop:bimorfisme-isomorfisme}.\end{indicacio}
\end{enumerate}

\section{Productes, equalitzadors i pullbacks}

\begin{enumerate}[resume]
  \item Demostreu que si \(p_A\circ u=p_A\circ v\) i \(p_B\circ u=p_B\circ v\), aleshores \(u=v\), i deduïu que \(\langle p_A,p_B\rangle=\id_{A\times B}\). (Model: secció~\ref{prac:productes}.)

  \item Demostreu que el producte de \(A\) i \(B\) és únic llevat d'isomorfisme, exhibint els isomorfismes inversos entre dues realitzacions qualssevol. \begin{indicacio}Seguiu el patró de la proposició~\ref{prop:producte-unic}: induïu \(\varphi,\psi\) per les propietats universals i useu la unicitat de factorització per veure \(\psi\circ\varphi=\id\).\end{indicacio}

  \item Demostreu que l'equalitzador \(e:E\rightarrowtail A\) de \(f,g:A\to B\) és un monomorfisme, i que qualsevol altre equalitzador dels mateixos \(f,g\) hi és equivalent com a part. (Model: secció~\ref{prac:equalitzadors-pullbacks}; enunciat: teorema~\ref{teo:equalitzador-universal}.)

  \item En un pullback de \(f:A\to C\) al llarg de \(g:B\to C\), demostreu que si \(f\) és un monomorfisme, la projecció oposada és un monomorfisme, i interpreteu-ho com que la imatge inversa d'una part és una part. \begin{indicacio}Vegeu la proposició~\ref{prop:pullback-mono}; cancel·leu \(f\) i useu la unicitat de la factorització del pullback.\end{indicacio}

  \item Siguin \(g:B\to A\) i una part \(m:S\rightarrowtail A\). Demostreu que \(y\in_X g^{-1}(m)\) si i només si \(g\circ y\in_X m\), i deduïu-ne que \(g^{-1}(\id_A)\equiv_B\id_B\) i que, si \(m\subseteq_A n\), aleshores \(g^{-1}(m)\subseteq_B g^{-1}(n)\). \begin{indicacio}La primera part és la proposició~\ref{prop:pertinenca-imatge-inversa}; per a les altres dues, compareu els elements generalitzats de les parts i useu la proposició~\ref{prop:inclusio-pertinenca}.\end{indicacio}

  \item Siguin \(h\colon C\to B\), \(g\colon B\to A\) i una part \(m\) de \(A\). Demostreu que \((g\circ h)^{-1}(m)\equiv_C h^{-1}\bigl(g^{-1}(m)\bigr)\). \begin{indicacio}Les dues parts tenen els mateixos elements generalitzats: apliqueu dues vegades la proposició~\ref{prop:pertinenca-imatge-inversa}. Alternativament, useu que dos quadrats cartesians encaixats formen un quadrat cartesià.\end{indicacio}

  \item \textbf{Axiomes usats.} Digueu quins axiomes nous del capítol (sumes, representació, productes) intervenen en cadascun d'aquests resultats, i quins es demostren només amb els del capítol~\ref{cap:aplicacions}: la composició de monomorfismes és un monomorfisme; tota secció és un monomorfisme; el testimoni d'una inclusió és únic; \(0,1\colon\mathbf1\to\mathbf2\) són monomorfismes; tot bimorfisme és un isomorfisme; l'equalitzador existeix; el pullback existeix. \begin{indicacio}Els tres primers només usen els axiomes estructurals i les definicions; el quart usa les sumes; el cinquè, la representació; els dos últims, productes i representació.\end{indicacio}
\end{enumerate}
